\documentclass[11pt,a4paper]{article}
\usepackage{fontspec}
\usepackage{xeCJK}
\usepackage{geometry}
\geometry{margin=25mm}
\setmainfont{Hiragino Sans}
\setCJKmainfont{Hiragino Sans}

\begin{document}
\begin{center}
{\Large 古物商許可申請書（内部提出用ひな形）}
\end{center}

\vspace{1em}
\noindent 申請 ID: APP-ANTIQUE-DEALER-001 \\
許認可種別: 古物商許可 \\
申請日: 令和8年7月14日

\section*{申請人}
\begin{tabular}{ll}
商号又は名称 & 株式会社MAL \\
本店所在地 & 〒102-0084 東京都千代田区二番町1 \\
\end{tabular}

\section*{代表者}
\begin{tabular}{ll}
氏名 & 段燕燕 \\
住所 & 〒102-0084 東京都千代田区二番町1 \\
\end{tabular}

\section*{事業所・対象}
\begin{tabular}{ll}
名称 & 株式会社MAL \\
所在地 & 〒102-0084 東京都千代田区二番町1 \\
管理者 & 段燕燕 \\
業態・備考 & 衣類・本・雑貨・時計・鞄等の日用品、PC・スマホ等の小物（買取は原則1点30万円以下） \\
\end{tabular}

\vfill
\small 公式参照: https://www.npa.go.jp/
\par\noindent\small 本ひな形は OpenOrgOS 内部整理用。提出前に管轄機関の最新様式を確認すること。
\end{document}
